% ---------------------------------------------------------------------------
% THC-LT-MP2: manuscript draft (theory and implementation only).
%
% Written for REVTeX 4-2 (AIP/JCP). If revtex4-2.cls is not installed the
% preamble falls back to `article' so that the file still compiles.
% ---------------------------------------------------------------------------
\IfFileExists{revtex4-2.cls}{%
  \documentclass[aip,jcp,amsmath,amssymb,preprint,superscriptaddress]{revtex4-2}%
  \newcommand{\FALLBACK}{0}%
}{%
  \documentclass[11pt]{article}%
  \newcommand{\FALLBACK}{1}%
}
\if 1\FALLBACK
  \usepackage[margin=1in]{geometry}
  \usepackage{amsmath,amssymb,graphicx}
  \newcommand{\affiliation}[1]{}
  \newcommand{\email}[1]{}
\fi

\usepackage{bm}
\usepackage{booktabs}

% ------------------------------------------------------------- shorthands --
\newcommand{\ket}[1]{\left|#1\right\rangle}
\newcommand{\bra}[1]{\left\langle#1\right|}
\newcommand{\Edir}{E_{\mathrm{dir}}}
\newcommand{\Eexc}{E_{\mathrm{exc}}}
\newcommand{\Emp}{E_{\mathrm{MP2}}}
\newcommand{\Naux}{N_{\mathrm{aux}}}
\newcommand{\Ngrid}{N_{\mathrm{grid}}}
\newcommand{\rr}{\mathbf{r}}
\newcommand{\re}{\operatorname{Re}}
\newcommand{\Ord}[1]{\mathcal{O}\!\left(#1\right)}

\begin{document}

\title{A tensor-hypercontraction, Laplace-transform formulation of
       second-order M{\o}ller--Plesset theory for two-component relativistic
       reference determinants with a bounded memory footprint}

\author{}
\affiliation{Department of Chemistry, Indian Institute of Technology Bombay,
             Powai, Mumbai 400076, India}

\date{\today}

\begin{abstract}
We describe a formulation and implementation of second-order
M{\o}ller--Plesset perturbation theory (MP2) for two-component,
spin--orbit-coupled reference determinants (exact two-component molecular
mean-field, X2CAMF) in which no four-index quantity is ever constructed. The
antisymmetrised correlation energy is separated into its direct (Coulomb) and
exchange Goldstone diagrams. The direct diagram is contracted through a
grid tensor-hypercontraction (THC) representation of the electron-repulsion
integrals combined with a Laplace transform of the doubles denominator; the
transform factorises the four orbital labels, and the diagram collapses onto
a sequence of $K\times K$ grid-space matrix multiplications, where $K$ is the
number of retained interpolation points. The exchange diagram, whose index
pattern the Laplace transform does not decouple and which is therefore
$\Ord{K^4}$ under THC, is instead evaluated by the resolution of the
identity with the occupied labels kept explicit; because those labels are
explicit no quadrature is required and the exact orbital-energy denominator
is used, so the exchange contribution carries no Laplace error. The
interpolation points and the THC core are built by a grid-blocked
interpolative separable density fitting procedure that never forms an
$\Ngrid\times\Ngrid$ or $\Ngrid\times n_{\mathrm{mo}}^2$ array and whose
least-squares normal equations are contracted over a single orbital or
atomic-orbital index rather than the orbital pair. The direct
kernel has an in-core path and an out-of-core path in which only two
$K\times K$ matrices are resident and the THC core is streamed in row
panels, so the peak memory is $\Ord{K^2}$ or $2K^2+\Ord{K\,n_{\mathrm{mo}}}$
rather than $\Ord{o^2v^2}$ or $\Ord{\Naux n_{\mathrm{mo}}^2}$. Spin-component
scaling of the two diagrams reproduces the spin-component-scaled and
scaled-opposite-spin variants; with the exchange diagram switched off the
method is a scaled-opposite-spin relativistic MP2 whose cost is $\Ord{K^3}$
per Laplace node. The implementation, its compiled kernels, and its
validation against a dense density-fitted MP2 in the same spinor molecular
orbital basis are described.
\end{abstract}

\maketitle

% ===========================================================================
\section{Introduction}
% ===========================================================================

Second-order M{\o}ller--Plesset perturbation theory is the least expensive
size-consistent wave-function correlation method and the entry point of the
correlation hierarchy. Its canonical cost is $\Ord{o^2v^2\Naux}$ in a
density-fitted (resolution-of-the-identity, RI) formulation,\cite{Whitten1973,
Dunlap1979,Vahtras1993,Feyereisen1993} with $o$ and $v$ the numbers of
occupied and virtual orbitals; the memory cost, if the doubles amplitudes or
the $(ov|ov)$ integral block are held, is $\Ord{o^2v^2}$. For a
two-component relativistic treatment on a spinor basis both counts refer to
complex arithmetic and to a spinor dimension roughly twice the number of
scalar atomic-orbital basis functions, so the $\Ord{o^2v^2}$ memory is the
binding constraint well before the arithmetic is: at a few thousand scalar
basis functions the four-index amplitude is measured in terabytes and the
three-index RI tensor in hundreds of gigabytes.

Two ideas reduce this. The Laplace transform of the energy
denominator,\cite{Almlof1991,Haser1992,Haser1993} $1/x=\int_0^\infty
e^{-xt}\,dt \approx \sum_g w_g e^{-x t_g}$, separates the four orbital
energies that enter an MP2 denominator into a product of per-orbital
exponential factors; those factors can be absorbed into the orbitals, after
which the direct (Coulomb) part of the correlation energy becomes a trace
over a product of two-index quantities.\cite{Jung2004,Nakajima2006} Tensor
hypercontraction\cite{Hohenstein2012,Parrish2012,Hohenstein2012b} represents
the electron-repulsion integral as a contraction of grid collocation
matrices $X_{Pp}$ (orbital $p$ evaluated at grid point $P$) with a small
core $Z_{PQ}$,
\begin{equation}
  (pq|rs) \;=\; \sum_{PQ}\rho_{pq}(P)\,Z_{PQ}\,\rho_{rs}(Q),
  \qquad
  \rho_{pq}(P)=\sum_{k}X^{*}_{kPp}X_{kPq},
  \label{eq:thc}
\end{equation}
so that any integral network can be rewritten in the $K$-dimensional grid
space. The interpolation points and the core are obtained by least squares
against the RI integrals; a robust and inexpensive route to the points is
interpolative separable density fitting (ISDF),\cite{Lu2015,Hu2017} which
selects them by a pivoted Cholesky (equivalently a column-pivoted QR)
factorisation of the collocation Gram matrix.

Combined, THC and the Laplace transform give an MP2 direct energy at
$\Ord{K^3}$ per Laplace node with only $\Ord{K^2}$ storage.\cite{Jung2004,
Song2016,Lee2020} The exchange energy is the difficulty: its integrand,
$(ia|jb)(ib|ja)$, ties the occupied label $i$ to the virtual labels of
\emph{both} factors, so no rearrangement of the Laplace product separates
the four grid sums, and a THC evaluation is irreducibly $\Ord{K^4}$ (or
$\Ord{o\,K^2 v}$ with the occupied labels retained), with a prefactor
$K/n_{\mathrm{mo}}\sim 10$ that makes it impractical past a thousand or so
functions. We therefore evaluate only the direct diagram by THC and Laplace,
and the exchange diagram by RI with the occupied labels kept explicit --- a
hybrid that keeps the memory bounded and, because the exchange labels are
explicit, keeps the exchange contribution free of quadrature error.

The target is the exact two-component molecular-mean-field (X2CAMF)
Hamiltonian,\cite{Liu2018,Zhang2022} for which BAGH already provides
Cholesky-decomposed spinor integrals and a THC representation used by its
coupled-cluster codes. The present method reuses that machinery. What
follows is the working equations (Sec.~\ref{sec:theory}) and the
implementation, including the grid-blocked ISDF build, the in-core and
out-of-core direct kernels, the screened exchange kernel, and frozen-core
handling (Sec.~\ref{sec:impl}); Sec.~\ref{sec:cost} collects the costs and
Sec.~\ref{sec:val} the numerical checks.

% ===========================================================================
\section{Theory}
\label{sec:theory}
% ===========================================================================

\subsection{The two diagrams}

In a spinor basis the MP2 correlation energy for a single-determinant
reference is
\begin{equation}
  \Emp = -\frac14 \sum_{ijab}
         \frac{\bigl|\bra{ij}\!\ket{ab}_{\!A}\bigr|^{2}}{D_{ijab}},
  \qquad
  \bra{ij}\!\ket{ab}_{\!A} = (ia|jb)-(ib|ja),
  \label{eq:emp2}
\end{equation}
with $D_{ijab}=\varepsilon_a+\varepsilon_b-\varepsilon_i-\varepsilon_j$,
indices $i,j$ occupied and $a,b$ virtual, $(ia|jb)$ the chemists'
electron-repulsion integral over spinors, and $\varepsilon$ the (semi-)
canonical spinor energies. Expanding the antisymmetriser,
\begin{equation}
  \bigl|(ia|jb)-(ib|ja)\bigr|^{2}
   = |(ia|jb)|^2 + |(ib|ja)|^2
     - 2\,\re\!\bigl[(ia|jb)(ib|ja)^{*}\bigr],
\end{equation}
and using that the first two terms are equal under
$i\!\leftrightarrow\!j,\ a\!\leftrightarrow\!b$ relabelling of the sum,
\begin{equation}
  \Emp = \Edir + \Eexc,
  \label{eq:split}
\end{equation}
\begin{align}
  \Edir &= -\frac12 \sum_{ijab} \frac{|(ia|jb)|^{2}}{D_{ijab}},
  \label{eq:edir}\\[2pt]
  \Eexc &= +\frac12 \sum_{ijab}
           \frac{\re\!\bigl[(ia|jb)(ib|ja)^{*}\bigr]}{D_{ijab}} .
  \label{eq:eexc}
\end{align}
Both sums have $D_{ijab}>0$ term by term. The exchange sum equals its own
complex conjugate under $i\!\leftrightarrow\!j$ together with the permutional
symmetry $(pq|rs)=(rs|pq)$, so $\Eexc$ is real and the $\re$ in
Eq.~\eqref{eq:eexc} may be taken inside.

\subsection{Direct diagram: THC and Laplace}

Write the Coulomb integral in the grid THC form Eq.~\eqref{eq:thc}, with the
collocation split into occupied and virtual blocks, $X^{o}_{kPi}$ and
$X^{v}_{kPa}$, and $k$ running over the two spinor components. Replace the
denominator by the Laplace quadrature,
$1/D_{ijab}=\sum_g w_g e^{-D_{ijab}t_g}$, valid because $D_{ijab}$ lies in a
bounded positive interval; the geometric (double-exponential in $1/h$)
quadrature of Ref.~\onlinecite{Takatsuka2008} on
$[2(\varepsilon_v^{\min}-\varepsilon_o^{\max}),\,
  2(\varepsilon_v^{\max}-\varepsilon_o^{\min})]$ is used, the factor two
being the doubles analogue of the factor three appropriate to triples.
The Boltzmann factor factorises as
$e^{-D_{ijab}t}=e^{(\varepsilon_i-\varepsilon_a)t}\,
e^{(\varepsilon_j-\varepsilon_b)t}$; splitting each half evenly and folding
it into the collocation,
\begin{equation}
  Y^{o}_{kPi}(t)=X^{o}_{kPi}\,e^{+\varepsilon_i t/2},
  \qquad
  Y^{v}_{kPa}(t)=X^{v}_{kPa}\,e^{-\varepsilon_a t/2},
  \label{eq:dress}
\end{equation}
so that the ``dressed'' pair density
$\tilde\rho_{ia}(P)=\sum_k Y^{o*}_{kPi}Y^{v}_{kPa}$ carries
$e^{(\varepsilon_i-\varepsilon_a)t/2}$. Substituting
Eqs.~\eqref{eq:thc},\eqref{eq:dress} into Eq.~\eqref{eq:edir} and summing the
orbital labels,
\begin{equation}
  \Edir = -\frac12 \sum_g w_g
          \sum_{PR} M_g[P,R]\,\bigl(Z\,M_g\,Z^{\dagger}\bigr)[P,R],
  \label{eq:edir-thc}
\end{equation}
\begin{align}
  M_g[P,R] &= \sum_{ia}\tilde\rho_{ia}(P)\,\tilde\rho_{ia}(R)^{*}
  \nonumber\\
  &= \sum_{kk'}
     \Bigl(\textstyle\sum_i Y^{o*}_{kPi}Y^{o}_{k'Ri}\Bigr)
     \Bigl(\textstyle\sum_a Y^{v}_{kPa}Y^{v*}_{k'Ra}\Bigr).
  \label{eq:Mg}
\end{align}
$M_g$ is a Hermitian $K\times K$ matrix, assembled as a Hadamard product of
two $K\times K$ Gram matrices for each of the four spinor-component pairs
$(k,k')$: eight complex matrix multiplications of shapes $(K\times o)(o\times
K)$ and $(K\times v)(v\times K)$. The bracketed quantity in
Eq.~\eqref{eq:edir-thc} is $N_g=Z M_g Z^{\dagger}$, two further $K\times K$
multiplications, and the node's contribution is the Frobenius inner product
$\sum_{PR}M_g[P,R]\,N_g[P,R]$, which is real and non-negative. The cost per
node is $\Ord{K^3}$; the largest array is $K\times K$.

\subsection{Exchange diagram: RI with exact denominators}

Because the Laplace product for the exchange integrand does not separate the
grid sums, we keep the occupied labels explicit and use the RI factorisation
$(ia|jb)=\sum_L L^{L}_{ia}L^{L}_{jb}$ that the reference already supplies
(the same three-index quantity from which the THC core is fitted). Define,
for an ordered occupied pair $(i,j)$,
\begin{equation}
  T^{ij}_{ab} = \sum_{L} L^{L}_{ia}\,L^{L}_{jb},
  \label{eq:Tij}
\end{equation}
a $v\times v$ matrix formed by one contraction over the auxiliary index.
Then Eq.~\eqref{eq:eexc} is
\begin{equation}
  \Eexc = \frac12 \sum_{i\le j} m_{ij}
          \sum_{ab}
          \frac{\re\!\bigl[T^{ij}_{ab}\,\bigl(T^{ij}_{ba}\bigr)^{*}\bigr]}
               {\varepsilon_a+\varepsilon_b-\varepsilon_i-\varepsilon_j},
  \label{eq:eexc-ri}
\end{equation}
with $m_{ij}=1$ for $i=j$ and $m_{ij}=2$ otherwise (the pair $(j,i)$
contributes the complex conjugate, with equal real part). No quadrature
enters Eq.~\eqref{eq:eexc-ri}: the denominator is the exact orbital-energy
gap, so $\Eexc$ is exact given the RI fit of the integrals. The cost is one
$\Ord{\Naux v^2}$ contraction and one $\Ord{v^2}$ weighted trace per pair,
$\Ord{o^2(\Naux+v)v^2}$ in total, with $\Ord{v^2}$ working memory.

For a spatially extended system with localised occupied
orbitals\cite{Pulay1983,Saebo1993} the summand of Eq.~\eqref{eq:eexc-ri}
decays with the separation of $i$ and $j$. A Cauchy--Schwarz bound,
\begin{equation}
  \bigl\|T^{ij}\bigr\|_{F}
   \;\le\; s_{ij}
   \;\equiv\; \sum_{L}\bigl\|L^{L}_{i\cdot}\bigr\|_{2}\,
                     \bigl\|L^{L}_{j\cdot}\bigr\|_{2},
  \label{eq:schwarz}
\end{equation}
is evaluated as a single $\Ord{o^2\Naux}$ matrix multiplication and used to
drop pairs with $s_{ij}$ below a threshold times $\max_{ij}s_{ij}$, taking
the pair count from $\Ord{o^2}$ towards $\Ord{o}$.

\subsection{Spin-component scaling}

Under spin--orbit coupling $S_z$ is not a good quantum number, so a
rigorous opposite-/same-spin partition of the MP2 energy does not exist.
Equations \eqref{eq:edir}--\eqref{eq:eexc} are the Coulomb and exchange
antisymmetrised diagrams; scaling them independently,
\begin{equation}
  \Emp = c_{\mathrm{os}}\,\Edir + c_{\mathrm{ss}}\,\Eexc,
  \label{eq:scs}
\end{equation}
recovers the non-relativistic spin-component-scaled\cite{Grimme2003} and
scaled-opposite-spin\cite{Jung2004} prescriptions in the weak-coupling
limit and provides their natural relativistic generalisation otherwise.
Setting $c_{\mathrm{ss}}=0$ removes the exchange diagram entirely; the
remaining method is a scaled-opposite-spin relativistic MP2 evaluated
purely by Eq.~\eqref{eq:edir-thc}, at $\Ord{K^3}$ per node and $\Ord{K^2}$
memory.

% ===========================================================================
\section{Implementation}
\label{sec:impl}
% ===========================================================================

The driver is a Python module; the two rate-limiting kernels ---
Eq.~\eqref{eq:edir-thc} and the pair loop of Eq.~\eqref{eq:eexc-ri} --- are
a compiled \textsc{pybind11} extension that calls level-3 BLAS
(\texttt{zgemm}) and is parallelised over Laplace nodes and grid panels, or
over occupied pairs, with OpenMP. Reference implementations of both
equations in dense array operations are retained and are used automatically
when the extension is absent; they are the object of the validation in
Sec.~\ref{sec:val}.

\subsection{ISDF factor construction}

The X2CAMF driver's own in-core THC factorisation forms an
$\Ord{\Ngrid n_{\mathrm{mo}}^{2}}$ array and is not used for this method;
the factors $X^{o}$, $X^{v}$, $Z$ are always built here, from the
density-fitted integrals, without any array of grid-squared or
grid$\times$orbital-squared size:
\begin{enumerate}
\item The two-component weighted collocation
  $X_{kPp}=w_P^{1/4}\sum_{\mu}\phi^{(k)}_{\mu}(\rr_P)\,C_{\mu p}$
  ($w_P$ the quadrature weight) is evaluated in blocks of grid points and
  held resident. It is $2\Ngrid n_{\mathrm{mo}}$ complex numbers --- a few
  hundred megabytes --- not the $\Ord{\Ngrid n_{\mathrm{mo}}^{2}}$ object,
  and every step below indexes into it.
\item Interpolation points are chosen by a pivoted Cholesky factorisation of
  the collocation Gram matrix
  $S_{PQ}=\sum_{pq}\rho_{P}(pq)^{*}\rho_{Q}(pq)
        =\sum_{kk'}\bigl|\textstyle\sum_p X_{kPp}X^{*}_{k'Qp}\bigr|^{2}$,
  the Schur-updated diagonal maintained explicitly. Each required column of
  $S$ is a few matrix--vector products into the resident collocation,
  $\Ord{\Ngrid n_{\mathrm{mo}}}$ work, not a re-evaluation of the atomic
  orbitals; no $\Ngrid\times\Ngrid$ matrix is stored. Selection stops when
  the updated diagonal falls below a relative floor, so $K$ is
  rank-adaptive.
\item The collocation is gathered on the $K$ retained points to give
  $X^{o}$ and $X^{v}$.
\item The THC core solves $Z = S^{-1} E\, S^{-1}$. Both $S$ and
  $E=\mathbf{T}\mathbf{T}^{\mathsf T}$ are formed by contracting a
  \emph{single} index, not the orbital pair $pq$:
  \begin{equation}
    S_{PQ}=\sum_{kk'}\bigl|(X_{k}X_{k'}^{\dagger})_{PQ}\bigr|^{2},
    \qquad
    \mathbf{T}_{PL}
      = \sum_{ab} L^{L}_{ab}
        \sum_{k,\,c\in\{\ell,s\}}
        \bigl(\bar X_{k}C_{c}^{\mathsf T}\bigr)_{\!Pa}
        \bigl(X_{k}\bar C_{c}^{\mathsf T}\bigr)_{\!Pb},
    \label{eq:lsthc}
  \end{equation}
  with $S$ assembled from four $(K\times n_{\mathrm{mo}})(n_{\mathrm{mo}}
  \times K)$ products and a Hadamard modulus-square, and the
  $\Naux\times n_{\mathrm{mo}}^{2}$ molecular-orbital fitting tensor
  $L^{L}_{pq}$ never built: $L^{L}_{ab}$ is the auxiliary Cholesky tensor
  (streamed over $L$), $C_{\ell},C_{s}$ the large/small-component MO
  coefficients, and $\bar{\phantom{x}}$ complex conjugation. The $(a,b)$
  block is formed in panels of $P$ and never exceeds
  $(p_{\mathrm{blk}}, n_{\mathrm{ao}}, n_{\mathrm{ao}})$. Since
  $n_{\mathrm{ao}}\approx n_{\mathrm{mo}}/2$ enters each sum in place of
  $n_{\mathrm{mo}}$, the solve costs roughly $n_{\mathrm{mo}}$ times fewer
  operations than the naive pair-index form for $S$ and about four times
  fewer for $E$. $S$ is regularised with a small multiple of its trace and
  inverted by a Cholesky solve (a pseudoinverse with an explicit cutoff as
  fallback); $Z$ is symmetrised.
\end{enumerate}
For large $K$ the core is written to disk and returned as a lazy dataset.
The block sizes are chosen from the available memory; the construction is
guarded so that a request whose $K\times K$ working set (a handful of such
matrices) would not fit is refused with a diagnostic rather than thrashing.

\subsection{Direct kernel: in-core and out-of-core}

The kernel loops over Laplace nodes. For each node it forms the dressed
collocation Eq.~\eqref{eq:dress}, builds $M_g$ from Eq.~\eqref{eq:Mg} as
eight \texttt{zgemm} calls and a Hadamard accumulation, and evaluates
Eq.~\eqref{eq:edir-thc}. Three residency regimes are selected automatically
from the estimated resident set against a memory budget:
\begin{itemize}
\item \emph{in-core}: $M_g$, $N_g$ and $Z$ are all held; the compiled
  kernel runs the entire node loop. Peak $\approx 4K^2$ complex plus the
  collocation.
\item \emph{out-of-core}: only $M_g$ and $W_g=M_g Z^{\dagger}$ are held and
  $Z$ is read in row panels from the on-disk dataset; $N_g=Z W_g$ and the
  inner product are accumulated panel by panel. Peak $\approx 2K^2$ complex
  plus the collocation, i.e.\ $2K^2+\Ord{K n_{\mathrm{mo}}}$.
\item a disk-backed $Z$ forces the out-of-core path.
\end{itemize}
The out-of-core path is pure BLAS-bound array algebra and is not moved into
the compiled extension; the in-core path is.

\subsection{Exchange kernel}

The Schwarz bound Eq.~\eqref{eq:schwarz} is formed from per-occupied
auxiliary-row norms and the surviving ordered pairs $(i\le j)$ are tiled by
blocks of the occupied index so that each three-index slice
$L^{L}_{i\in\mathcal{I},a}$ that must be resident is bounded. For each tile
the compiled kernel loops over its pairs; per pair it forms $T^{ij}$
(Eq.~\eqref{eq:Tij}) by one \texttt{zgemm} contracting the auxiliary index
of the two strided slices, then accumulates
$m_{ij}\sum_{ab}\re[T^{ij}_{ab}(T^{ij}_{ba})^{*}]/
       (\varepsilon_a+\varepsilon_b-\varepsilon_i-\varepsilon_j)$
into a thread-private scalar. The dressing of Eq.~\eqref{eq:dress} is
\emph{not} applied here --- there is no Laplace grid --- and no dressed copy
of the three-index integral is made.

\subsection{Frozen core}

The standard frozen-core keywords select a number of lowest occupied
spinors to exclude. The occupied collocation $X^{o}$, the occupied energy
vector, the three-index integral $L_{ia}$, and the interval on which the
Laplace nodes are placed are all sliced to the active occupied space before
either diagram is evaluated; the virtual space and $K$ are unchanged. The
frozen core acts identically on both diagrams, and the correlation energy
reproduces a dense frozen-core RI-MP2 in the same active space to the THC
and quadrature accuracy.

% ===========================================================================
\section{Computational cost and memory}
\label{sec:cost}
% ===========================================================================

\begin{table}[t]
\caption{Formal cost and peak working memory. $o$, $v$ occupied and virtual
spinor counts; $n_{\mathrm{ao}}\approx n_{\mathrm{mo}}/2$ scalar atomic
orbitals; $\Naux$ auxiliary functions; $K$ retained interpolation points
($K\sim c\,n_{\mathrm{mo}}$, $c\sim 8$--$15$); $\Ngrid$ parent grid points;
$\Ngrid^{\mathrm{LT}}$ Laplace nodes; $p$ the $P$-block size in the ISDF
build.}
\label{tab:cost}
\begin{ruledtabular}
\begin{tabular}{lll}
step & cost & peak memory \\
\colrule
ISDF collocation \& points & $\Ngrid\,K\,n_{\mathrm{mo}}$ &
   $\Ngrid\,n_{\mathrm{mo}} + \Ngrid\,K$ \\
ISDF core, $S$ & $K^{2}\,n_{\mathrm{mo}}$ & $K^{2}$ \\
ISDF core, $E$ & $K\,\Naux\,n_{\mathrm{ao}}^{2}$ &
   $K^{2}+p\,n_{\mathrm{ao}}^{2}$ \\
direct, in-core     & $K^{3}\,\Ngrid^{\mathrm{LT}}$ & $4K^{2}$ \\
direct, out-of-core & $K^{3}\,\Ngrid^{\mathrm{LT}}$ &
   $2K^{2}+K\,n_{\mathrm{mo}}$ \\
exchange (RI)       & $o^{2}(\Naux+v)\,v^{2}$ & $v^{2}$ \\
\end{tabular}
\end{ruledtabular}
\end{table}

Table~\ref{tab:cost} collects the counts. No step scales as $o^2v^2$ or
$\Naux n_{\mathrm{mo}}^2$ in memory, and the single-index form of
Eq.~\eqref{eq:lsthc} keeps the LS-THC solve off the pair index as well.
The direct diagram is the rate-determining arithmetic at $\Ord{K^3}$ per
node; with the geometric Laplace quadrature $\Ngrid^{\mathrm{LT}}\sim
10$--$20$ suffices for $10^{-6}$--$10^{-7}$ relative accuracy on the
denominator in a valence basis. An all-electron basis, whose tight
correlating virtuals push the interval endpoint $\varepsilon_v^{\max}$
very high, inflates the geometric rule to $\sim 80$ nodes; a minimax
quadrature would restore the small count and is a natural refinement. The
exchange
diagram, at the cost of a conventional RI-MP2 exchange, is the practical
ceiling on the full calculation: its $v\times v$ per-pair work is not yet
truncated to pair domains, so it is affordable to $v\sim 3000$--$4000$.
Beyond that, $c_{\mathrm{ss}}=0$ removes it and the direct term alone scales
to the largest spinor bases with a memory footprint of a few $K^2$ complex
numbers held or streamed.

% ===========================================================================
\section{Validation}
\label{sec:val}
% ===========================================================================

The working equations were checked at three levels.

\emph{Kernel identities.} On random collocation factors and a random
Hermitian core, the compiled direct kernel reproduces the reference
evaluation of Eq.~\eqref{eq:edir-thc} to machine precision
($\sim 10^{-15}$ relative), and both reproduce the defining sum
$-\tfrac12\sum_{ijab}|(ia|jb)|^2\sum_g w_g e^{-D_{ijab}t_g}$ built densely
from the same factors; the residual against the exact $1/D$ evaluation
falls with the Laplace spacing as expected. The out-of-core direct path,
with the core handed back in row panels, agrees with the in-core path to
$\sim 10^{-14}$. The compiled exchange kernel, summed over the full pair
list, reproduces the exact-denominator reference of Eq.~\eqref{eq:eexc-ri}
to $\sim 10^{-15}$ and the Laplace form of the same quantity to the
quadrature accuracy.

\emph{Against a dense density-fitted MP2.} For a closed-shell atom
(neon, X2CAMF, correlation-consistent double-$\zeta$ with a matching
auxiliary set) the three-index integrals were transformed into the X2CAMF
spinor molecular-orbital basis and the antisymmetrised MP2 energy formed
densely from Eq.~\eqref{eq:emp2}. The interpolation-separable factors were
then built by the procedure of Sec.~\ref{sec:impl} and the kernel run: the
exchange diagram matched the dense value to $\sim 10^{-15}$; the direct
diagram and the total matched to the tensor-hypercontraction fit and
Laplace quadrature error, $\sim 10^{-6}$--$10^{-7}$ in the correlation
energy at the default grid and spacing. A frozen-core run
($2$ frozen spinors) reproduced a dense frozen-core RI-MP2 in the same
active space to the same accuracy.

\emph{Independent paths.} The in-core and out-of-core direct evaluations,
the compiled and reference kernels, and the scaled-opposite-spin
($c_{\mathrm{ss}}=0$) path were checked pairwise for consistency on the
same inputs.

\emph{A production-scale run.} As a size and stability check the method was
applied to $n$-hexane in the all-electron Dyall AE2Z basis on a
two-component X2CAMF reference --- $648$ spinor molecular orbitals, $50$
occupied of which $12$ are frozen as core --- on a single shared
$46\,\mathrm{GB}$ node. The chunked ISDF procedure held a
$0.22\,\mathrm{GB}$ collocation, selected $K=4909$ interpolation points
from a $10\,640$-point parent grid, and ran the LS-THC solve of
Eq.~\eqref{eq:lsthc} with a peak working set of a few $K\times K$ matrices;
the direct and exchange kernels then returned
$\Edir=-1.4399$, $\Eexc=+0.5057$ and $\Emp=-0.9342\,E_{\mathrm h}$. The
single-index reformulation of the normal equations cut the kernel wall time
by about one third relative to the pair-index form, at no change in the
energy beyond floating-point reassociation ($\sim 10^{-8}\,E_{\mathrm h}$),
and the run never approached the node's memory. As anticipated above, the
all-electron basis drove the geometric Laplace rule to $86$ nodes for the
direct term.

% ===========================================================================
\section{Conclusions}
% ===========================================================================

A memory-bounded MP2 for two-component relativistic references follows from
treating the two Goldstone diagrams by different factorisations: the direct
diagram by tensor hypercontraction and a Laplace transform, which reduce it
to $K\times K$ grid-space matrix multiplications at $\Ord{K^3}$ per node
with $\Ord{K^2}$ storage; the exchange diagram, which the Laplace transform
does not decouple, by the resolution of the identity with the occupied
labels explicit, so that it needs no quadrature and is exact to the RI fit.
The interpolation factors are built by a grid-blocked ISDF procedure that
avoids every grid-squared array and whose least-squares normal equations
are contracted over a single index, the direct kernel has an out-of-core
path that streams the THC core, and the exchange kernel takes a Schwarz
pair screen. A twenty-atom all-electron test completes on a single
$46\,\mathrm{GB}$ node. With the exchange diagram removed the method is a
scaled-opposite-spin relativistic MP2 that scales to the largest spinor
bases; the remaining refinements are a minimax Laplace quadrature to
contain the node count with all-electron bases and a pair-domain
truncation of the exchange, both left for future work.

\appendix

\section{Grid THC / ISDF conventions}

The spinor pair density on the parent grid is
$\rho_{pq}(P)=\sum_{k=1}^{2}X^{*}_{kPp}X_{kPq}$ with
$X_{kPp}=w_P^{1/4}\,\varphi^{(k)}_{p}(\rr_P)$ and $\varphi^{(k)}_p$ the
$k$-th component of spinor $p$. The least-squares THC core minimises
$\sum_{pqrs}\bigl|(pq|rs)-\sum_{PQ}\rho_{pq}(P)Z_{PQ}\rho_{rs}(Q)\bigr|^2$
with $(pq|rs)$ taken from the RI fit; its stationary condition is
$S Z S = E$ with $S_{PQ}=\sum_{pq}\rho_{pq}(P)^{*}\rho_{pq}(Q)$ and
$E_{PQ}=\sum_{pqrs}\rho_{pq}(P)(pq|rs)\rho_{rs}(Q)$, matching the
convention already used by the X2CAMF THC coupled-cluster code so that the
factors are interchangeable between methods.

\begin{thebibliography}{99}

\bibitem{Whitten1973} J. L. Whitten,
  \emph{Coulombic potential energy integrals and approximations},
  J. Chem. Phys. \textbf{58}, 4496 (1973).

\bibitem{Dunlap1979} B. I. Dunlap, J. W. D. Connolly and J. R. Sabin,
  \emph{On some approximations in applications of X$\alpha$ theory},
  J. Chem. Phys. \textbf{71}, 3396 (1979).

\bibitem{Vahtras1993} O. Vahtras, J. Alml{\"o}f and M. W. Feyereisen,
  \emph{Integral approximations for LCAO-SCF calculations},
  Chem. Phys. Lett. \textbf{213}, 514 (1993).

\bibitem{Feyereisen1993} M. Feyereisen, G. Fitzgerald and A. Komornicki,
  \emph{Use of approximate integrals in ab initio theory. An application in
  MP2 energy calculations}, Chem. Phys. Lett. \textbf{208}, 359 (1993).

\bibitem{Almlof1991} J. Alml{\"o}f,
  \emph{Elimination of energy denominators in M{\o}ller--Plesset
  perturbation theory by a Laplace transform approach},
  Chem. Phys. Lett. \textbf{181}, 319 (1991).

\bibitem{Haser1992} M. H{\"a}ser and J. Alml{\"o}f,
  \emph{Laplace transform techniques in M{\o}ller--Plesset perturbation
  theory}, J. Chem. Phys. \textbf{96}, 489 (1992).

\bibitem{Haser1993} M. H{\"a}ser,
  \emph{M{\o}ller--Plesset (MP2) perturbation theory for large molecules},
  Theor. Chim. Acta \textbf{87}, 147 (1993).

\bibitem{Jung2004} Y. Jung, R. C. Lochan, A. D. Dutoi and M. Head-Gordon,
  \emph{Scaled opposite-spin second order M{\o}ller--Plesset correlation
  energy: An economical electronic structure method},
  J. Chem. Phys. \textbf{121}, 9793 (2004).

\bibitem{Nakajima2006} T. Nakajima and K. Hirao,
  \emph{The Douglas--Kroll--Hess approach}, and Laplace-transformed MP2
  developments therein, Chem. Rev. (representative);
  see also A. Takatsuka, S. Ten-no and W. Hackbusch,
  J. Chem. Phys. \textbf{129}, 044112 (2008).

\bibitem{Takatsuka2008} A. Takatsuka, S. Ten-no and W. Hackbusch,
  \emph{Minimax approximation for the decomposition of energy denominators
  in Laplace-transformed M{\o}ller--Plesset perturbation theories},
  J. Chem. Phys. \textbf{129}, 044112 (2008).

\bibitem{Hohenstein2012} E. G. Hohenstein, R. M. Parrish and T. J. Mart{\'i}nez,
  \emph{Tensor hypercontraction density fitting. I. Quartic scaling second-
  and third-order M{\o}ller--Plesset perturbation theory},
  J. Chem. Phys. \textbf{137}, 044103 (2012).

\bibitem{Parrish2012} R. M. Parrish, E. G. Hohenstein, T. J. Mart{\'i}nez and
  C. D. Sherrill,
  \emph{Tensor hypercontraction. II. Least-squares renormalization},
  J. Chem. Phys. \textbf{137}, 224106 (2012).

\bibitem{Hohenstein2012b} E. G. Hohenstein, R. M. Parrish, C. D. Sherrill and
  T. J. Mart{\'i}nez,
  \emph{Communication: Tensor hypercontraction. III. Least-squares tensor
  hypercontraction for the determination of correlated wavefunctions},
  J. Chem. Phys. \textbf{137}, 221101 (2012).

\bibitem{Lu2015} J. Lu and L. Ying,
  \emph{Compression of the electron repulsion integral tensor in tensor
  hypercontraction format with cubic scaling cost},
  J. Comput. Phys. \textbf{302}, 329 (2015).

\bibitem{Hu2017} W. Hu, L. Lin and C. Yang,
  \emph{Interpolative separable density fitting decomposition for accelerating
  hybrid density functional calculations},
  J. Chem. Theory Comput. \textbf{13}, 5420 (2017).

\bibitem{Song2016} C. Song and T. J. Mart{\'i}nez,
  \emph{Atomic orbital-based SOS-MP2 with tensor hypercontraction},
  J. Chem. Phys. \textbf{144}, 174111 (2016).

\bibitem{Lee2020} J. Lee, L. Lin and M. Head-Gordon,
  \emph{Systematically improvable tensor hypercontraction: Interpolative
  separable density fitting for molecules},
  J. Chem. Theory Comput. \textbf{16}, 243 (2020).

\bibitem{Grimme2003} S. Grimme,
  \emph{Improved second-order M{\o}ller--Plesset perturbation theory by
  separate scaling of parallel- and antiparallel-spin pair correlation
  energies}, J. Chem. Phys. \textbf{118}, 9095 (2003).

\bibitem{Pulay1983} P. Pulay,
  \emph{Localizability of dynamic electron correlation},
  Chem. Phys. Lett. \textbf{100}, 151 (1983).

\bibitem{Saebo1993} S. Saeb{\o} and P. Pulay,
  \emph{Local treatment of electron correlation},
  Annu. Rev. Phys. Chem. \textbf{44}, 213 (1993).

\bibitem{Liu2018} J. Liu and L. Cheng,
  \emph{An atomic mean-field spin--orbit approach within exact two-component
  theory for a non-perturbative treatment of spin--orbit coupling},
  J. Chem. Phys. \textbf{148}, 144108 (2018).

\bibitem{Zhang2022} C. Zhang and L. Cheng,
  \emph{Atomic mean-field approach within exact two-component theory},
  representative application, J. Phys. Chem. A \textbf{126}, 4537 (2022).

\end{thebibliography}

\end{document}
